Large language models produce poorly calibrated confidence estimates, expressing high certainty on questions they answer incorrectly. We investigate whether chain-of-thought (CoT) prompting improves calibration by enabling models to ``self-read'' their own uncertainty signals. Through systematic mechanism verification on TruthfulQA, we demonstrate that CoT prompting produces reasoning chains containing epistemic hedging markers (82\% presence rate, mean 2.84 markers per output) that negatively correlate with verbalized confidence (Spearman $r=-0.315$, $p<10^{-16}$, $n=664$). All four steps of the hypothesized causal mechanism---reasoning generation (100\% rate), hedging marker presence, structural positioning (100\% compliance), and confidence adjustment---are empirically validated. These findings support a self-reading interpretation: models integrate in-context uncertainty signals from their own reasoning into confidence judgments, providing a zero-shot prompting approach to improved calibration without model fine-tuning.
